\documentclass[11pt]{article}

\usepackage[margin=1in]{geometry}
\usepackage{amsmath,amssymb,amsthm,mathtools}
\usepackage{enumitem}
\usepackage{hyperref}
\usepackage[nameinlink,capitalize]{cleveref}

\hypersetup{colorlinks=true,linkcolor=blue,citecolor=blue,urlcolor=blue}

\newtheorem{theorem}{Theorem}
\newtheorem{proposition}{Proposition}
\newtheorem{lemma}{Lemma}
\newtheorem{corollary}{Corollary}
\newtheorem{definition}{Definition}
\newtheorem{remark}{Remark}
\newtheorem{example}{Example}

\newcommand{\Z}{\mathbb Z}
\newcommand{\F}{\mathbb F}
\newcommand{\E}{\mathbb E}
\newcommand{\Var}{\operatorname{Var}}
\newcommand{\rank}{\operatorname{rank}}
\newcommand{\Span}{\operatorname{Span}}
\newcommand{\one}{\mathbf 1}
\newcommand{\C}{\mathbb C}

\title{Kummer-Fourier Collision Variance and Weighted Hyperplane Matroid Sums}
\author{Jared Wilder}
\date{2026-05-13}

\begin{document}
\maketitle

\begin{abstract}
This note corrects a tempting but false linear product formula for Kummer collision excess.  The corrected identity is exact finite algebra: for a CRT product of local Kummer obstruction indicators, the collision variance of the induced obstruction-vector distribution is a Plancherel sum whose local factors are squared complex characteristic functions.  Equivalently, the variance is a weighted hyperplane-arrangement sum over the dual of the Kummer log matroid.  The formula is unconditional and finite; no analytic equidistribution theorem, PHNC input, or Erd\H{o}s--Graham resolution is asserted.
\end{abstract}

\section{Setup}

Let \(\mathcal Q\) be a finite set of primes and put
\[
P=\prod_{q\in\mathcal Q}q.
\]
Let \(\ell\) be prime, and let \(x_q\in\F_\ell^r\) be the Kummer log row attached to \(q\) for a fixed rank-\(r\) base vector \(\mathbf a=(a_1,\ldots,a_r)\).  Write
\[
\Gamma_q=\langle a_1,\ldots,a_r\rangle\subset\F_q^*
\]
and
\[
p_q=\frac{|\Gamma_q|}{q-1}.
\]
For \(m\in(\Z/P\Z)^*\), define the local obstruction indicator
\[
B_q(m)=\one_{\Gamma_q}(-m_q^{-1}),
\]
where \(m_q\) denotes the image of \(m\) in \(\F_q^*\).  By the Chinese remainder theorem, as \(m\) varies uniformly over \((\Z/P\Z)^*\), the variables \(B_q\) are independent Bernoulli variables with
\[
\mathbb P(B_q=1)=p_q.
\]

\begin{definition}[Bernoulli-Kummer obstruction vector]
Define
\[
V(m)=\sum_{q\in\mathcal Q} B_q(m)x_q\in\F_\ell^r.
\]
Let
\[
N(v)=\#\{m\in(\Z/P\Z)^*:V(m)=v\},
\qquad v\in\F_\ell^r,
\]
and
\[
\bar N=\frac{\varphi(P)}{\ell^r}.
\]
The collision excess of \(N\) is
\[
\mathcal C_\ell(N)=\sum_{v\in\F_\ell^r}(N(v)-\bar N)^2.
\]
\end{definition}

\section{The false linear formula}

The following expression is a useful failed route:
\[
\ell^{-2r}\varphi(P)^2
\sum_{A\in\F_\ell^r\setminus\{0\}}
\left(
\prod_{q\mid P}
\left[
1+\frac{\ell\one_{A\cdot x_q=0}-1}{q-1}
\right]
-1
\right).
\]
It is not equal to \(\mathcal C_\ell(N)\) in general.  The failure is structural: Plancherel requires squared complex Fourier coefficients, while the displayed product is a linear real proxy.  The factor \(\ell^{-2r}\) and the internal subtraction are also incompatible with the standard finite Fourier normalization.

\begin{example}[Minimal sign and normalization failure]
Let \(r=1\), \(\ell=2\), \(\mathcal Q=\{q\}\), and \(x_q=1\).  Then \(V(m)=B_q(m)\in\F_2\).  Hence
\[
N(1)=|\Gamma_q|,
\qquad
N(0)=q-1-|\Gamma_q|.
\]
Therefore
\[
\mathcal C_2(N)
=
2\left(|\Gamma_q|-\frac{q-1}{2}\right)^2
=
\frac{(q-1)^2}{2}(1-2p_q)^2.
\]
The false linear formula gives a different value.  For instance, at \(q=3\), \(a_1=2\), one has \(p_q=1\), so \(\mathcal C_2(N)=2\), while the false expression gives a nonmatching linear value.  The identity is false as stated.
\end{example}

\section{Exact Kummer-Fourier collision variance}

Let
\[
\zeta_\ell=e^{2\pi i/\ell}.
\]
For \(A\in\F_\ell^r\), define the finite Fourier coefficient
\[
\widehat N(A)=\sum_{v\in\F_\ell^r}N(v)\zeta_\ell^{A\cdot v}.
\]

\begin{theorem}[Exact Kummer-Fourier collision variance]\label{T:kummer-fourier-collision}
With notation as above,
\[
\mathcal{C}_\ell(N)=
\frac{\varphi(P)^2}{\ell^r}
\sum_{A\in\F_\ell^r\setminus\{0\}}
\prod_{q\in\mathcal Q}
\left[
1-2p_q(1-p_q)
\left(1-\cos\frac{2\pi(A\cdot x_q)}{\ell}\right)
\right].
\]
Here \(A\cdot x_q\) is represented by any integer lift modulo \(\ell\); the cosine is well-defined.
\end{theorem}

\begin{proof}
By finite Fourier inversion and Plancherel on \(\F_\ell^r\),
\[
\mathcal C_\ell(N)=
\ell^{-r}\sum_{A\ne0}|\widehat N(A)|^2.
\]
The \(A=0\) term is excluded because subtracting \(\bar N\) removes the constant component.

For \(A\in\F_\ell^r\),
\[
\widehat N(A)
=
\sum_{m\in(\Z/P\Z)^*}\zeta_\ell^{A\cdot V(m)}
=
\sum_{m\in(\Z/P\Z)^*}
\prod_{q\in\mathcal Q}\zeta_\ell^{B_q(m)(A\cdot x_q)}.
\]
Under the uniform measure on \((\Z/P\Z)^*\), CRT makes the variables \(B_q\) independent, so
\[
\widehat N(A)
=
\varphi(P)
\prod_{q\in\mathcal Q}
\E\left[\zeta_\ell^{B_q(A\cdot x_q)}\right].
\]
For a fixed \(q\), since \(B_q\) is Bernoulli with parameter \(p_q\),
\[
\E\left[\zeta_\ell^{B_q(A\cdot x_q)}\right]
=
(1-p_q)+p_q\zeta_\ell^{A\cdot x_q}.
\]
Thus
\[
|\widehat N(A)|^2
=
\varphi(P)^2
\prod_{q\in\mathcal Q}
\left|(1-p_q)+p_q\zeta_\ell^{A\cdot x_q}\right|^2.
\]
Writing \(\theta_q=2\pi(A\cdot x_q)/\ell\),
\[
\left|(1-p_q)+p_q e^{i\theta_q}\right|^2
=
(1-p_q+p_q\cos\theta_q)^2+(p_q\sin\theta_q)^2
=
1-2p_q(1-p_q)(1-\cos\theta_q).
\]
Substituting into Plancherel gives the formula.
\end{proof}

\begin{corollary}[Binary specialization]\label{C:binary-specialization}
When \(\ell=2\),
\[
\mathcal{C}_2(N)=
\frac{\varphi(P)^2}{2^r}
\sum_{A\in\F_2^r\setminus\{0\}}
\prod_{q\in\mathcal Q}
\left[
1-4p_q(1-p_q)\bigl(1-\one_{A\cdot x_q=0}\bigr)
\right].
\]
\end{corollary}

\begin{proof}
For \(\ell=2\), \(\cos(\pi t)=1\) when \(t=0\) in \(\F_2\), and \(\cos(\pi t)=-1\) when \(t=1\).  Substituting in \cref{T:kummer-fourier-collision} gives the claim.
\end{proof}

\section{Matroid interpretation}

Let \(M_\ell(\mathcal Q;\mathbf a)\) be the representable Kummer log matroid on ground set \(\mathcal Q\), represented by the rows \(x_q\in\F_\ell^r\).  For \(A\ne0\), the condition \(A\cdot x_q=0\) says that the row \(x_q\) lies in the hyperplane
\[
H_A=\{x\in\F_\ell^r:A\cdot x=0\}.
\]
Thus \cref{T:kummer-fourier-collision} is a weighted sum over hyperplanes of the dual arrangement.

\begin{definition}[Kummer hyperplane weight]
For \(A\ne0\), define
\[
W_A(\mathcal Q)=
\prod_{q\in\mathcal Q}
\left[
1-2p_q(1-p_q)
\left(1-\cos\frac{2\pi(A\cdot x_q)}{\ell}\right)
\right].
\]
\end{definition}

\begin{corollary}[Weighted hyperplane-arrangement form]
The collision excess is
\[
\mathcal C_\ell(N)
=
\frac{\varphi(P)^2}{\ell^r}
\sum_{A\ne0}W_A(\mathcal Q).
\]
For \(\ell=2\), \(W_A\) depends only on the subset of rows outside the hyperplane \(H_A\):
\[
W_A(\mathcal Q)=
\prod_{q:x_q\notin H_A}
\left[1-4p_q(1-p_q)\right].
\]
\end{corollary}

\begin{remark}[Not an ordinary Tutte evaluation]
The formula is not, in general, a single ordinary Tutte-polynomial evaluation \(T_M(x,y)\).  The weights are prime-dependent through \(p_q\), and for \(\ell>2\) they depend on the nonzero value of \(A\cdot x_q\), not only on whether it vanishes.  The honest description is a weighted hyperplane-arrangement sum, equivalent in the binary uniform-weight case to a multivariate Tutte/Potts-style specialization.
\end{remark}

\section{Dependency status}

\begin{itemize}[leftmargin=2em]
\item The theorem is exact finite algebra: dependency class \(E\).
\item The proof uses only CRT independence and finite Fourier Plancherel.
\item No PHNC, BVK, Kummer-BV, GRH, Chebotarev error term, or Erd\H{o}s--Graham conclusion is used.
\item The result strengthens the finite algebraic stack by replacing a false linear product with the correct squared-modulus product.
\end{itemize}

\end{document}
